\scanpage{4}
\mspage{3}

Condition plus forte : \deleted{\unsure{la topologie de $I$ est quotient pour}}
\inserted{(pour $\varphi$) l'application (continue \uline{surjective}) $\varphi$ est \uline{ouverte}}
\deleted{\unsure{celle de}}
\[
\deleted{\unsure{X_\varphi\text{ est fidèle}}}\quad\Longleftrightarrow\quad
\varphi^{-1}(\overline J)=\overline{\varphi^{-1}(J)}
\quad\Longleftrightarrow\quad
\varphi^{-1}(\overline j)=\overline{\varphi^{-1}(j)}\quad\text{i.e.}
\]
\[
\begin{array}{cc}
\Vert&\Vert\\
\displaystyle\bigcup_{j\in\overline J}X_j^*&
\displaystyle\overline{\bigcup_{j\in J}X_j^*}=\bigcup_{j\in J}\overline{X_j^*}
\end{array}
\]
\[
\overline{X_j^*}=\bigcup_{i\in\overline j}X_i^*\ (\overset{\mathrm{déf}}=X_j).
\]

Conditions pour une pseudo-partition $(X_i^*)_{i\in I}$ \inserted{\uline{loc. finie}} indexée par un ens
\uline{préordonné} $I$, i.e. $\varphi:X\longrightarrow I$ application, en supposant $\varphi$ \uline{continue}
\deleted{que} les $X_i=\bigcup_{j\leq i}X_j^*$ fermés (donc les $X_i^*$ loc. fermés)

\unsure{2°|4°}) Top. de $I$ est finale pour $\varphi$ (i.e. quotient si $\varphi$ surjective) $\Longleftrightarrow$
\[
(X_J^*\text{ fermé}\Rightarrow J\text{ fermé})\Longleftrightarrow
(X_L^*\text{ ouvert}\Rightarrow L\text{ ouvert})
\]
\margin{si $\varphi$ surj. i.e. 1° $\Uparrow$}

3°) $\unsure{\overline{X_j^*}|\overline{X_i^*}}=\bigcup_{j\leq i}X_j^*$ pour tt $\unsure{j|i}\in I$ i.e. $\varphi$ est ouverte

1°) $\varphi$ surj i.e. les $X_i^*\ne\varnothing$

\margin{\uline{NB} Moyennant 1°, 3°, on
récupère l'ordre de $I$ comme
la rel\unsure{ation} d'inclusion entre les
$\overline{X_i^*}$ \ldots}

\unsure{La donnée} $(X_i^*)_{i\in I}$ satisfaisant 1°) et 3°) équivaut à
\unsure{une} : une relation d'équiv. \uline{loc\unsure{t.} finie} et \uline{ouverte}
\deleted{\unsure{telle que les classes soient}}
et $I$ est l'espace quotient (avec la relation de \inserted{pré}ordre
de spécialisation), \deleted{\illegible{3}} l'ens $I$ est \uline{ordonné} ssi
\[
(\forall i,j\quad\overline{X_i^*}=\overline{X_j^*}\Rightarrow i=j)
\Longleftrightarrow(\text{les }X_i\text{ sont loc. fermés}).
\]

\margin{cas $I$ \uline{ordonné} $\Updownarrow$}
\uline{Stratification} \uline{\unsure{stricte}}

famille loc. finie \inserted{\uline{croissante}} $(X_i)_{i\in I}$ de fermés, telle que $\forall j<i$
implique que $X_j$ rare dans $X_i$, et $\forall i,j\in I$,
on ait
\[
X_i\cap X_j=\bigcup_{k\in\overline i\cap\overline j}X_k.
\]

Dans le cadre de la top. modérée, il faut voir \unsure{comment|au moins}
\inserted{pour} toute partition par des parties loc. fermées \uline{modérées},
il existe une plus petite \inserted{stratification en des parties}
\deleted{\unsure{partition finie}} \unsure{telle} qui
\deleted{\unsure{engendre}} la partition donnée.

\mspage{4}
\margin{\diagram{\[\begin{array}{cc|cc}
\subset\Rightarrow\times&\times\Rightarrow\subset&\subset&\subseteq\\
\unsure{\cap}&\uparrow&&\\
&\subset&&
\end{array}\]}}

$(X_i^*)_{i\in I}$ une \inserted{partition indexée} \deleted{famille} \inserted{de $X$}
\deleted{(loc. finie de parties disjointes de $X$,
de réunion $X$, donc \unsure{une} $X\xrightarrow{\varphi}I$)}. Cond. équiv.

a) $\forall i\in I$, \deleted{les $X_i^*\ne\varnothing$ (i.e. $\varphi$ surj, i.e. une partition), et}
$\overline{X_i^*}$ est saturé \deleted{\unsure{pour les $X_i^*$, loc. fermés}}

b) $\exists\leq$, une \deleted{relation} \inserted{pré}ordre sur $I$, telle que l'on ait $\forall i\in I$
\[
\overline{X_i^*}=\bigcup_{j\leq i}X_j^*
\]
(cette \deleted{relation} préordre est unique, \ldots
\[
i\leq j\Longleftrightarrow\overline{X_i^*}\subset\overline{X_j^*}
\quad(\Longleftrightarrow X_i^*\subset\overline{X_j^*})\ )
\]
c) La top. quotient de $I$ est telle que $\varphi:X\longrightarrow I$ soit ouverte,

i.e. $\forall U\in\operatorname{Ouv}(X)$, le saturé est ouvert

Sous ces conditions, la top. quotient de $I$ est bien
associée à \deleted{\unsure{sa relation d'ordre}} préordre, qui est aussi
la relation de spécialisation pour cette topologie.
Pour que $I$ soit ordonné, il f. et s. que les
$X_i^*$ soient loct. fermés. On dit alors que la partition est
\deleted{\unsure{stratifiée}} \unsure{(i.e. les $X_i^*$ saturés)}.
\inserted{\unsure{loc. finie}, pièces loc. fermées, relation d'équiv. ouverte}

On a une \uuline{correspondance} 1--1 entre les partitions
\inserted{\uline{stratifiantes}} de $X$ de cette forme, i.e. les relations d'équiv. loc.
finies et ouvertes, \inserted{sur $X$, à fibres loc. fermées}, et les familles (à iso près)
$(X_i)_{i\in I}$ indexées par un ens. d'indices $I$ \uline{ordonné}
\inserted{\uline{loc. fini}},
qui sont telles que 1°) familles \uline{croissantes} de parties
\inserted{\uline{fermées}} \deleted{\unsure{croissantes}} de $X$, \uline{recouvrant} $X$, et 2°) $\forall i,j\in I$
\[
X_i\cap X_j=\bigcup_{k\in\overline i\cap\overline j}X_k
\]
4°) \deleted{$\unsure{X_\alpha}$} $\forall i,j\in I$, $j<i$
$\Longrightarrow$ \inserted{$X_j$ rare dans $X_i$}
\deleted{\unsure{(telle que $X_i^*=X_i$)}}

3°) $\forall i$, $X_i^*\ne\varnothing$
\inserted{``déf'' $X_i\setminus\bigcup_{j<i}X_j$}

\margin{\uline{NB} Moyennant \unsure{1°}, 2°, 4°, 3° équivaut à
\unsure{la cond. plus faible}
3° bis $(\forall i\in I)$ $\unsure{X_i}\ne\varnothing$.}

Moyennant 1°) et 2°), en posant $X_i^*=X_i\setminus\bigcup_{j<i}X_j^*$,
les $X_i^*$ forment une ps. partition \inserted{\uline{loc. finie}} de $X$ par des
ens. loc. fermés, celle-ci satisfait les conditions
d'une \uline{partition} \uline{stratifiée} \uline{stricte} \inserted{(avec l'ordre donné de $I$)} ssi
on a 3°) et 4°).

\uline{Stratification} : Conditions \unsure{2°) et 1°)}. \uline{Strat. stricte} : ``fidèle'' (i.e.
3°) et condition de densité (4°) : \unsure{Partitions stratifiées} $\Longleftrightarrow$ \unsure{strat. strictes}

\margin{3}
